\section*{September 7, 2026: Preserve address-geocoding evidence}
\textit{Agent-drafted reconstruction; location acceptance remains under review.}

Jacob requested a required checkpoint before the full rerun: return to all 30
unresolved initial predecessor requests, comprising 24 with multiple candidate
polygons and six with none. Each must be traced through the later construction
pipeline and recorded as resolved, ineligible under an established rule, or needing
judgment. A successful build does not complete this review. The checkpoint is
recorded in the construction-cleaning README, with the request identifiers and
status filters needed to recover the cases.

Address geocoding can now be replayed from complete, preserved service responses.
The old exports saved selected matches but omitted the other candidates. September 7
queries to the Census and Chicago services returned full responses for 196 distinct
historical addresses. Applying the original rules reproduces all fields of both
old 211-row chosen-match tables, including coordinates. Request selection and
candidate acceptance now run locally in construction cleaning; a separate
acquisition task preserves the responses and exact queries. Ordinary builds verify
the pinned source hashes and make no geocoder calls.

This exact replay also preserves a weakness in the old Census rule. It requires a
unique match, matching house number, and coordinates within Chicago's broad bounds,
but does not compare the street. Four accepted responses differ from the requested
address: 1236 N Troy Street becomes S Troy Street; 763 W 15th Place becomes Street;
and 3609 and 3623 W 50th Street become Place. The old reference-point calculation
used three of these Census responses; a City point match supplied the correct
requested street for 3623. All four projects survive in the frozen paper input.
Of the three using Census, one lies within 500 feet in that input. This establishes
a reason to investigate, not the correct physical location: a historical source
address can itself be wrong. No effect on estimates has been calculated.

The proposed general rule requires agreement on house number, direction, street
name, and street type after standard abbreviation normalization, while allowing an
omitted trailing unit label. Jacob's decision remains pending. The restored code
therefore retains the old rules as a comparison baseline. This step does not
approve their matches for the full rerun or introduce density exclusions.

The current request set has 216 rows. Its original 211 requests are unchanged;
five new requests lack a selected historical address. Census still accepts 194
requests and Chicago 114 under the old rules. The later parcel lookup also needed
17 additional year/parcel queries: six were available in the initial source, and
11 newly queried pairs returned six polygons. Combining sources preserves all
11,794 original preferred-source polygons, including their exact geometry, and
does not refresh previously queried empty results. Current preferred coverage has
15,691 component-year requests, 17 more than before.

\textit{Reproduction and limits.} Work follows \path{2ca07cf} on
\path{sales_sample_exploration}. Run the address-request and two geocode report
targets in \path{tasks/new_construction_cleaning/code/}. Source snapshots and
download recipes belong to \path{download_construction_address_geocodes} and
\path{download_construction_historical_parcels}. The original-coverage comparison
uses preserved address history and current exact-PIN coordinate availability;
it reproduces the original request set exactly. Missing queries, duplicate keys,
API errors, empty requests, no-address and no-match responses, independent report
recovery, and unchanged builds were checked. Later physical-location review and
the full pipeline remain unfinished. Shared Make rules now check upstream report
targets as well as data outputs; tests verified missing-report recovery and
changed-input propagation from document, ordinary-task, and nested-task paths.
The paper input is unchanged; estimates were
not rerun. Logbook compilation retains earlier holds on unfinished audit products
and previously verified development documents. Earlier commercial-footprint,
ground-up-evidence, and commercial-ledger reports are also held at their recorded
baselines: refreshing those downstream stages stops at preferred predecessor
points outside the pinned spatial-query scope. The new geocoding reports build
through their actual dependencies without those holds.
